%-------------------------
% Resume in Latex
% Author : Sourabh Bajaj
% Website: https://github.com/sb2nov/resume
% License : MIT
%------------------------

\documentclass[letterpaper,11pt]{article}

\usepackage{latexsym}
\usepackage[empty]{fullpage}
\usepackage{titlesec}
\usepackage{marvosym}
\usepackage[usenames,dvipsnames]{color}
\usepackage{verbatim}
\usepackage{enumitem}
\usepackage[pdftex]{hyperref}
\usepackage{fancyhdr}
\usepackage{fontawesome5}


\pagestyle{fancy}
\fancyhf{} % clear all header and footer fields
\fancyfoot{}
\renewcommand{\headrulewidth}{0pt}
\renewcommand{\footrulewidth}{0pt}

% Adjust margins
\addtolength{\oddsidemargin}{-0.5in}
\addtolength{\evensidemargin}{-0.375in}
\addtolength{\textwidth}{1in}
\addtolength{\topmargin}{-.5in}
\addtolength{\textheight}{1.0in}

\urlstyle{same}

\raggedbottom
\raggedright
\setlength{\tabcolsep}{0in}

% Sections formatting
\titleformat{\section}{
  \vspace{-5pt}\scshape\raggedright\Large
}{}{0em}{}[\color{black}\titlerule \vspace{-5pt}]

%-------------------------
% Custom commands
\newcommand{\resumeItem}[2]{
  \item\small{
    \textbf{#1}{: #2 \vspace{-3pt}}
  }
}

\newcommand{\subItem}[1]{
  \item\small{
    {#1 \vspace{-2pt}}
  }
}

\newcommand{\resumeSubheading}[4]{
  \vspace{0pt}\item
    \begin{tabular*}{0.97\textwidth}{l@{\extracolsep{\fill}}r}\vspace{-2pt}
      \textbf{#1} & #2 \\\vspace{-2pt}
      \textit{\small#3} & \textit{\small #4} \\
    \end{tabular*}\vspace{-6pt}
}

\newcommand{\courseSubheading}[1]{
  \vspace{0pt}\item
    \begin{tabular*}{0.97\textwidth}{l@{\extracolsep{\fill}}r}
        \textbf{#1}
        \end{tabular*}\vspace{-6pt}
}

\newcommand{\resumeSubItem}[2]{\resumeItem{#1}{#2}\vspace{-4pt}}

\renewcommand{\labelitemii}{$\circ$}

\newcommand{\resumeSubHeadingListStart}{\begin{itemize}[label={}, leftmargin=0em]}
\newcommand{\resumeSubHeadingListEnd}{\end{itemize}}
\newcommand{\resumeItemListStart}{\begin{itemize}[leftmargin=1.5em]}
\newcommand{\resumeItemListEnd}{\end{itemize}\vspace{-6pt}}

%-------------------------------------------
%%%%%%  CV STARTS HERE  %%%%%%%%%%%%%%%%%%%%%%%%%%%%


\begin{document}

%----------HEADING-----------------

\begin{center}
    \textbf{\LARGE \href{https://brandonho667.github.io/}{Brandon Ho}} \\ \vspace{5pt}
    \small \faPhone* {(858) 330-9711} \hspace{1pt} $|$
    \hspace{1pt} \href{mailto:brandonho667@gmail.com} {\faEnvelope \hspace{2pt} {brandonho667@gmail.com}} \hspace{1pt} $|$
    \hspace{1pt} \href{https://github.com/brandonho667}{\faGithub \hspace{2pt} {brandonho667}} \hspace{1pt} $|$
    \hspace{1pt} \href{https://www.linkedin.com/in/bbho/}{\faLinkedin \hspace{2pt} {bbho}}
    \\ \vspace{-3pt}
\end{center}

%-----------EDUCATION-----------------
\section{Education}
  \resumeSubHeadingListStart
    \resumeSubheading
      {University of California San Diego}{GPA: 3.93/4 (M.S.)}
      {B.S. Computer Engineering (Magna Cum Laude), M.S. Intelligent Systems, Robotics \& Control}{June 2024}
        \resumeItemListStart
            \resumeItem{Relevant Coursework}{Software Engineering, Deep Learning, Operating Systems, Computer Vision, Autonomous Vehicles, Sensing \& Estimation, Planning \& Learning, Robot Sys Design \& Implementation, Statistical Learning}
        \resumeItemListEnd


  \resumeSubHeadingListEnd


%-----------EXPERIENCE-----------------
\section{Experience}
  \resumeSubHeadingListStart
  \resumeSubheading
      {Relativity Space}{Long Beach, CA}
      {Robotics Software Engineer II}{September 2024 - Current}
      \resumeItemListStart
        \subItem
          {Building a flexible, \textbf{real-time} robotics software platform (\textbf{C++}) to support large metal additive manufacturing}
        \resumeItem{Project Management}
          {Successfully executed \textbf{1-year R\&D project} to revamp core SW middleware with team of 3 engineers, wrote and iterated on roadmap, presented monthly to stakeholders, delegated work via \textbf{Jira}}
        \resumeItem{Experimental Design}
        {Designed and executed \textbf{6+ weld trials} characterizing weld param to bead geometry relationships; consulted weld engineers, documented and standardized procedures, logged all telemetry to cloud for full traceability, performed \textbf{Python} data analysis and visualization to communicate results to simulation team}
        \resumeItem{RT Middleware}
          {Built novel RT-capable \textbf{Zenoh+EVL}-based RMW for in-band/out-of-band comms; migrated platform from shmem/DDS, increasing throughput \textbf{10x}, reliability \textbf{1.5x}, reducing codebase redundancy \textbf{15\%}}
        \resumeItem{RT Linux Kernel}
          {Automated \textbf{Xenomai4} build/packaging via \textbf{Ansible}; contributed upstream patches fixing \href{https://gitlab.com/Xenomai/xenomai4/linux-evl/-/commit/e8fe20b2467a9ed75d58efc1d3afb50b032993fa}{network stack} packet leak and \href{https://gitlab.com/Xenomai/xenomai4/linux-evl/-/commit/fe63838c8a0c1c249006fb41c735c4904947d2d6}{page faults} in real-time processes, \textbf{reducing core SW crashes by 90\%}}
        \resumeItem{Sensor Integration}
      {Integrated Gocator 2D laser profilometer with real-time \textbf{RANSAC circle fitting} to extract weld bead geometry at \textbf{10Hz}, increasing part geometry \textbf{data completeness 20x}, reducing noise by \textbf{30\%}}
        \resumeItem{Closed-loop Control}
        {Implemented \textbf{PID} w/ anti-windup for wire feed speed, \textbf{reducing bead variation by 30\%}}
        \resumeItem{Web UI}
        {Developed \textbf{TypeScript/React} command center with \textbf{Foxglove WebSocket}, migrated from rosbridge reducing crashes by \textbf{50\%}; collaborated with weld engineers to build UI features improving weld process control}
      \resumeItemListEnd
  \resumeSubheading
      {Relativity Space}{Long Beach, CA}
      {Robotics Software Engineer Intern}{June 2023 - September 2023}
      \resumeItemListStart
        \subItem
          {Exposing \textbf{OPC-UA} interface on Fronius welders to real-time context for Wire Arc Additive Manufacturing}
        \resumeItem{ROS2 Driver}
          {Wrote an OPC-UA driver in \textbf{ROS2 (C++)} to read and write over RT (shared memory) and non-RT contexts, exposing \textbf{40+ new parameters} and increasing controllable welding parameters by \textbf{5x}}
        \resumeItem{System Design}
          {Refactored welder state machine and inverted object dependencies, \textbf{reducing 1000+ lines of code}}
        \resumeItem{UI Dev}
          {Created web interface for weld engineers to interact with OPC-UA parameters in \textbf{TypeScript}}
      \resumeItemListEnd
    % \resumeSubheading
    %   {Garmin}{Brea, CA}
    %   {Software Engineer Intern}{June 2022 - December 2022}
    %   \resumeItemListStart
    %     \subItem
    %       {Automating release processes and mock radio UI layer on the Tuner Team (Garmin AutoOEM Radio)}
    %     \resumeItem{Process Automation}
    %       {Utilized Python REST APIs to automate page generation with user-configurable metadata}
    %     \resumeItem{UI Automation}
    %       {Automated mock UI creation and deployment using MQTT \& Node-RED for \textbf{214 D-Bus services}}
    %   \resumeItemListEnd
    % \resumeSubheading
    %   {Yonder Dynamics}{La Jolla, CA}
    %   {Software Lead}{October 2020 - Present}
    %   \resumeItemListStart
    %     \subItem
    %       {Lead software dev for rover to autonomously path and perform tasks in the University Rover Competition}
    %     \resumeItem{Project Management}
    %     {Leading a team of \textbf{20+ SWE w/ SCRUM} and spearheading documentation on Notion}
    %     \resumeItem{Autonomous Exploration}
    %       {Implementing AR-tag detection with computer vision and researching RRT* search algorithm to update a live AR-tag probability map in C++, \textbf{increasing rover exploration speed by \texttt{$\sim$}10x}}
    %     \resumeItem{Sensor Fusion}
    %     {Fusing GPS with local position from IMU and stereo camera for a \textbf{$\sim$20\% increase in accuracy}}
    %     \resumeItem{Obstacle Avoidance}{Created a \textbf{ROS} obstacle detection and avoidance pipeline in C++ with \textbf{RealSense} stereo camera to efficiently convert 3D pointclouds at 30Hz into occupancy grid and output avoidance waypoints}
    %   \resumeItemListEnd
    \resumeSubheading
      {InflammaSense}{La Jolla, CA}
      {Software Developer}{January 2021 - September 2022}
      \resumeItemListStart
        \subItem
          {Built medical device to record and analyze neural signals for early signs of sepsis, published \textbf{Scientific Reports} (\href{https://www.nature.com/articles/s41598-022-21817-w}{link})}
        \resumeItem{Data Pipeline}
          {Designed a robust async data pipeline with the \textbf{ZeroMQ} Python API for the wireless display, transmission, storage (\textbf{PostgreSQL}), analysis of neural data for an array of \textbf{100+ medical devices (\texttt{$\sim$}155 Mbps)}}
        \resumeItem{Project Management}
          {Documented a comprehensive system model and \textbf{leading a team of 5+ SWE} with Asana}
        \resumeItem{UI Dev}
          {Built a GUI to display high throughput EKG data (8 kHz) w/ \textbf{PyQt}, enabling real-time data monitoring}
      \resumeItemListEnd

    % \resumeSubheading
    %   {Bae Lab}{UCSD Nanoengineering}
    %   {Undergraduate Research Assistant}{September 2020 - Present}
    %   \resumeItemListStart
    %     \subItem
    %       {Researching computational methods to control hydrogel mechanics via analysis of microscopy images and modulus}
    %   \resumeItemListEnd

    % \resumeSubheading
    %   {Coding for Impact}{Your Life Counts}
    %   {Web Developer}{July 2020 - Present}
    %   \resumeItemListStart
    %     \subItem
    %       {Assisting suicide prevention nonprofit Your Life Counts in web development, reaching out to 1300+ people}
    %     \resumeItem{WordPress}
    %       {Using PHP/CSS/HTML/JS to create and design lightweight dynamic plug-ins (video carousel)}
    %   \resumeItemListEnd

  \resumeSubHeadingListEnd


%-----------PROJECTS-----------------
\section{Projects}
  \resumeSubHeadingListStart
    \resumeSubheading
        {Autonomous Rover Path Planning} {Yonder Dynamics}
        {Software Lead}{September 2020 - June 2023}
        \resumeItemListStart
        \subItem
          {Lead software dev for rover to autonomously path and perform tasks in the University Rover Competition}
        \resumeItem{Autonomous Exploration}
          {Implemented RRT* search in C++, \textbf{increasing rover exploration speed \texttt{$\sim$}10x}}
        \resumeItem{Obstacle Avoidance}{Created a \textbf{ROS} obstacle detection and avoidance pipeline in C++ with \textbf{RealSense} stereo camera to efficiently convert 3D pointclouds at 30Hz into occupancy grid and output avoidance waypoints}
        \resumeItemListEnd
    % \resumeSubheading
    %     {SEMPro}{Bae Lab}
    %     {Using Explainable AI to Learn Structure-Property Insights from Hydrogel Microscopy Images}{Spring 2021}
    %     \resumeItemListStart
    %     \resumeItem{Web Scraping}
    %       {Developed NLP scraping algorithm (Selenium+BeautifulSoup), compiling data from \textbf{1100+ articles}}
    %     \resumeItem{Deep Learning}
    %       {Optimized and trained a ResNext model in Pytorch to predict hydrogel modulus from dataset of microscopy images, able to predict within a range of 1 log Pa with \textbf{an accuracy of $\sim$90\%}}
    %     \resumeItem{Explainable AI}
    %         {Implemented regression activation mapping on a CNN to expose learned structure-property insights}
    %     \resumeItemListEnd
    % \resumeSubheading
    %     {\href{https://github.com/brandonho667/autow}{Autow}} {Autonomous Vehicles}
    %     {Autonomous tow hitching using Computer Vision}{Fall 2022}
    %     \resumeItemListStart
    %     \subItem
    %         {Led a team of two mechanical and one software engineer in programming an RC car to autonomously hitch to a trailer}
    %     \resumeItem{ROS2}
    %         {Built an async callback system with \textbf{ROS2} for seamless cohesion between driver and autonomous control}
    %     \resumeItem{Computer Vision}
    %         {Utilized \textbf{OpenCV} for localization w/ ArUco tags, resulting in a \textbf{90\% hitching accuracy}}
    %     \resumeItemListEnd

    % \resumeSubheading
    %     {Vehicle Trajectory Prediction}
    %     {Automatic tow hitching algorithm using Computer Vision}{Fall 2022}
    %     \resumeItemListStart
    %     \subItem
    %         {Built and trained deep learning models for class-wide Kaggle competition, placing \textbf{top 10 across 50 teams}}
    %     \resumeItem{Data Analysis}
    %         {Analyzed and preprocessed trajectory metrics, \textbf{improving baseline performance by 3x}}
    %     \resumeItem{LSTM}
    %         {Implemented an LSTM encoder-decoder model w/ attention \textbf{improving model performance by 11x}}
    %     \resumeItemListEnd

    % \resumeSubheading
    %     {UCSD Connect}{IEEE Quarterly Projects}
    %     {Connecting UCSD students during the COVID-19 social isolation crisis}{December 2020}
    %     \resumeItemListStart
    %     \subItem
    %         {Led a team of 5 developers to create a college social media app and win \textbf{2nd Place} in the 2020 UCSD IEEE QP}
    %     \resumeItem{Full-stack App Development}
    %         {Utilized Android Studio to design app UI and back-end (login, matching, chat)}
    %     \resumeItem{System Management}
    %         {Directed communication between front/back-end and workflow \& issue creation in GitHub}
    %     \resumeItemListEnd

    % \resumeSubheading
    %     {PersonAI Chat}{UCSD Summer Program for Incoming Students}
    %     {AI Chatbot with Customizable Personas}{August 2020}
    %     \resumeItemListStart
    %     \subItem{Led a team of 4 developers to create the back-end for the AI chat site, ProfessorUWU}
    %     \resumeItem{TensorFlow/OpenAI GPT}
    %         {Implemented TensorFlow with OpenAI GPT's Transformer-based language model on custom dialogue datasets to create unique personas to chat with}
    %     \resumeItemListEnd

  \resumeSubHeadingListEnd

%--------SKILLS------------
\section{Skills}
 \resumeSubHeadingListStart
   \vspace{-2pt}\item
        \textbf{Programming Languages:}{ Python, C++, Java, SQL, JS, PHP, HTML/CSS, TypeScript }
   \vspace{-5pt}

   \vspace{-2pt}\item
        \textbf{Technologies:}{ AWS, Git, Jupyter Notebook, Linux, Docker, Confluence, Jira, Gerrit, Gitlab CI, Claude }
    \vspace{-5pt}

    \vspace{-2pt}\item
        \textbf{Frameworks:}{ ROS2, OPC-UA, PyTorch, MQTT, Zenoh, Xenomai, NumPy, Pandas, SciKit, Ansible }
    \vspace{-5pt}
 \resumeSubHeadingListEnd


%-------------------------------------------
\end{document}
